% !TEX program = pdflatex
\documentclass[11pt,letterpaper]{article}

\usepackage[T1]{fontenc}
\usepackage[utf8]{inputenc}
\usepackage[margin=0.9in,headheight=14pt]{geometry}
\usepackage{microtype}
\usepackage[dvipsnames]{xcolor}
\usepackage{enumitem}
\usepackage{titlesec}
\usepackage{fancyhdr}
\usepackage[round,authoryear]{natbib}
\usepackage{url}
\usepackage{hyperref}

\definecolor{WarnerBlue}{HTML}{005A9C}
\definecolor{LinkBlue}{HTML}{005F9E}
\definecolor{SoftGray}{HTML}{555555}

\hypersetup{
  colorlinks=true,
  linkcolor=LinkBlue,
  citecolor=LinkBlue,
  urlcolor=LinkBlue,
  pdftitle={Module 2 Reflection Journal: Planning for Access, Agency, and Justice in Invisible Invaders},
  pdfauthor={Piter Z. Garcia Bautista},
  pdfsubject={EDU486 Module 2 Reflection},
  pdfkeywords={justice-centered STEM, neurodivergence, chronic illness, science modeling, environmental justice, camp pedagogy}
}

\titleformat{\section}
  {\large\bfseries\color{WarnerBlue}}
  {}{0pt}{}
\titlespacing*{\section}{0pt}{1.25em}{0.45em}

\setlist[itemize]{leftmargin=1.4em,itemsep=0.35em,topsep=0.4em}
\setlength{\parindent}{0pt}
\setlength{\parskip}{0.72em}
\renewcommand{\baselinestretch}{1.08}
\widowpenalty=10000
\clubpenalty=10000
\displaywidowpenalty=10000

\pagestyle{fancy}
\fancyhf{}
\lhead{\small EDU486 Module 2 Journal}
\rhead{\small Access, Agency, and Justice}
\cfoot{\small\thepage}
\renewcommand{\headrulewidth}{0.4pt}

\newcommand{\keyclaim}[1]{\textbf{#1}}

\begin{document}

\hypersetup{pageanchor=false}
\begin{titlepage}
  \thispagestyle{empty}
  \vspace*{1.15in}

  {\centering
    {\Huge\bfseries\color{WarnerBlue} Module 2 Reflection Journal\par}
    \vspace{0.28in}
    {\LARGE Planning for Access, Agency, and Justice\par}
    \vspace{0.08in}
    {\LARGE in Invisible Invaders\par}

    \vfill

    {\large Piter Z. Garcia Bautista\par}
    \vspace{0.12in}
    {\large EDU486: Integrating Science and Technology\par}
    {\large Warner School of Education\par}
    \vspace{0.22in}
    {\large July 22, 2026\par}
  }

  \vfill

  \noindent\color{SoftGray}\small
  \textbf{Journal prompt.} After reviewing the Heather Table of course readings, describe the principles that inform camp pedagogy, make reading-to-principle connections explicit when possible, and draw on other course experiences.
\end{titlepage}
\hypersetup{pageanchor=true}
\pagenumbering{arabic}

\section{From a List of Activities to a Pedagogy}

At the beginning of Module 2, I could describe many things I wanted youth to do: collect sand, compare samples, inspect suspected particles, build a model, discuss injustice, and create art for advocacy. What I could not yet explain as clearly was the pedagogy connecting those activities. Reviewing my Heather Table alongside the Module 2 readings helped me see that a good camp is not a crowded schedule of engaging tasks. It is a learning system in which youth can understand why the next investigation matters, recognize that their knowledge belongs in science, revise ideas with evidence, and choose how they will act.

\keyclaim{My central principle is that inclusion and rigor are the same design problem.} Both require me to make the intellectual work visible, gather better evidence of what youth are thinking, and revise the environment before I interpret a learner's silence, fatigue, movement, communication style, or need for time as a lack of ability. This connects directly to my Puzzle Plan and EQUITAS Access-to-Agency work with neurodivergent, chronically ill, disabled, multilingual, culturally and linguistically diverse, and multiply marginalized learners. Module 2 gave that commitment a more concrete instructional form.

For our Invisible Invaders camp, this means organizing the week around the anchoring phenomenon, \emph{``Humans aren't trash cans, but we have plastics in our bodies,''} and the driving question, \emph{``How does plastic break into tiny pieces, move through the environment and into living bodies, and how can we stop it fairly?''} The phenomenon invites investigation, but the way we facilitate that investigation will determine whether youth experience science as something they help build or as another set of conclusions delivered to them.

\section{Principle 1: Coherence Must Be Visible to Learners}

\citet{reiser2021} argue that a storyline should be coherent from the students' perspective: youth questions and emerging explanations should create a felt need for the next investigation. \citet{cherbow2026} press this point further by showing that a well-designed curriculum cannot guarantee a coherent experience. Teachers still have to build the relationship among students, curricular resources, and the evolving storyline. Backward design helps me identify the final understanding and acceptable evidence before choosing activities \citep{wiggins2005},\footnote{The Module 2 schedule labels this reading as 2021; the assigned course PDF is the 2005 expanded second edition cited here.} but a backward plan remains only my plan unless youth can see and influence the path.

\keyclaim{For Invisible Invaders, coherence is now something youth can see, not something hidden in our planning document.} We will keep the driving question, the day's purpose, the complete activity sequence, and our Four Gotta-Haves visible. The four ideas are that Earth systems are interconnected, water transports materials, living organisms respond to environmental conditions, and human infrastructure influences natural systems. Each connecting question attaches to one of those ideas so that an activity earns its place by helping youth explain the phenomenon.

This principle changed after meeting the Freedom Scholars and staff on July 21. Youth asked for short explanations, one direction at a time, and a chance to begin doing quickly. That feedback did not ask us to simplify the science. It asked us to remove the executive-function burden of holding an invisible sequence in working memory. A visual agenda, a completed starting example, and brief demonstrations make the causal problem easier to enter while leaving the thinking demanding. When a learner joins late, takes a break, or needs to re-enter after pain, sensory overload, or attention loss, the storyline remains available without requiring a public explanation.

\section{Principle 2: Models Are Public, Revisable Thinking}

The modeling toolkit helped me move beyond treating a model as a polished poster made at the end of a unit. \citet{windschitl2013} frame models as public representations that make both observable events and unobservable mechanisms available for critique and revision. Science and literacy also strengthen one another when learners observe, investigate, talk, read, write, and communicate for an authentic purpose rather than completing disconnected language exercises \citep{ippolito2018}. Local litter inquiries show how firsthand evidence can lead toward explanation and public communication \citep{smetana2019,wisbrock2020}.

\keyclaim{Our camp model must preserve change, uncertainty, and evidence limits.} On Field Friday, youth will compare four controlled jars and examine traceable Charlotte Beach sand samples. On Monday, they will build a first source-to-body model. During the week, they will add evidence about transport, friction and wear, living systems, infrastructure, and unequal decision-making. On Thursday, they will revise a shared model rather than replace the first one with an adult answer. Friday's KEEP, QUESTION, CHANGE routine will make one final evidence-supported revision or justify why an existing relationship should remain.

This also requires honest claim language. A visually suspected particle is not chemically confirmed plastic. Finding microplastics in a body or food source establishes exposure, but it does not by itself prove a particular route or health effect. Youth can label claims as \emph{supported}, \emph{possible}, or \emph{still unknown}. Those labels are not signs that the investigation failed. They are part of learning how science protects people from false certainty and how responsible advocacy can remain strong without exaggeration.

Multiple participation routes belong inside this same scientific practice. A learner may contribute by speaking, pointing, drawing arrows, writing a label, dictating to a partner, taking an evidence photograph, or asking to pass and return. These are not easier versions of modeling. Each route still asks the learner to connect evidence, mechanism, and claim. The representation changes; the intellectual demand does not.

\section{Principle 3: Place, Community, and Family Knowledge Are Science}

\citet{seiki2021} describe border crossing between home and school by positioning familial knowledge alongside formal science rather than asking learners to leave one world behind. \citet{starr2026} similarly argue that outdoor learning is not just science moved to a different location. Place changes relationships among youth, families, educators, histories, and socioecological systems. The transdisciplinary river course described by \citet{nesmith2021} shows why environmental problems cannot be understood through science content alone; policy, economics, ethics, art, history, and community knowledge shape both the problem and possible action.

\keyclaim{Charlotte Beach and the places youth already know are not scenery for our lesson; they are part of the evidence and the question of responsibility.} We will return to Field Friday observations throughout the week, map how water and infrastructure connect places, and ask who makes decisions, who benefits, who carries burdens, and whose knowledge is usually recognized. Youth may connect the model to a neighborhood, bus route, laundry practice, food system, family story, or object they know. We will not collect private addresses or use ZIP codes as a shortcut for identity or risk.

I also need to treat connection as an invitation, not a demand for disclosure. A question about plastic in human bodies can intersect with illness, disability, food insecurity, environmental racism, family work, or medical trauma. No learner should have to disclose a diagnosis, body history, or family experience to prove that the topic matters. Youth can work through an invented scenario, a public place, a provided image, or a community-scale example. That choice protects identity safety while keeping community knowledge available as an intellectual resource.

\section{Principle 4: Identity, Recognition, and Joy Belong Inside Rigor}

\citet{mercier2021} show how identi-beads and identi-badges can help youth recognize multiple forms of STEM participation, including investigating, inventing, designing, tinkering, conservation, and care for others. The important move is not the bead itself. It is making a wider range of contributions visible to oneself and to the group. \citet{muhammad2023} brings identity, skills, intellect, criticality, and joy into one curricular frame, reminding me that justice-centered learning cannot stop after naming harm.

\keyclaim{I want every Freedom Scholar to leave with evidence that their way of contributing mattered to the science community we built.} We will notice careful observation, revision, connection-making, advocacy, creativity, care, and the questions that redirect the group. We will deliberately recognize quiet and behind-the-scenes contributions, not only the fastest speaker, the student who handles equipment first, or the person most comfortable presenting publicly. Recognition cannot become a ranking system in which some identities are treated as more scientific than others. Youth identities also remain fluid; a bead or role names a contribution from that moment, not a permanent category.

Joy, in this camp, is not a prize added after rigorous work. It appears when youth experience curiosity, movement, humor, friendly challenge, beauty in advocacy art, meaningful relationships, and the power to imagine a more just future. July 21 feedback made this practical: include movement and social connection, but also a seated or lower-energy path; show an example without requiring imitation; and give early finishers a meaningful extension instead of more repetitive work. Joy becomes compatible with chronic illness and neurodivergence only when participation does not require performing endless energy.

\section{Principle 5: Access Must Be Built Into the Thinking Environment}

\citet{liljedahl2020} challenged me to examine classroom routines that produce compliance without much thinking. I value the emphasis on rich tasks, visible collaborative work, carefully chosen hints and extensions, and teachers answering questions in ways that return agency to learners. PICRAT adds a related technology question: What relationship does the learner have with the tool, and does the tool replace, amplify, or transform the practice \citep{kimmons2020}? UDL reinforces planning for multiple ways to engage, perceive information, and express understanding from the beginning \citep{cast2024}.

\keyclaim{However, no single posture, communication mode, group arrangement, or technology can stand in for thinking.} A vertical whiteboard may make ideas public for one learner and create pain, fatigue, motor, sensory, or social barriers for another. Random grouping may widen collaboration, but it can also remove a stabilizing relationship. Verbal questioning may reveal reasoning, or it may measure processing speed and speech access. My responsibility is to preserve the goal of visible, collaborative thinking while offering standing, seated, tabletop, visual, verbal, partner, photo, and low-energy routes.

This is where my prior work most clearly presses against a mechanical use of a framework. I have repeatedly seen access barriers interpreted as thinking barriers, especially for neurodivergent and chronically ill learners whose knowledge is not expressed through the expected pace, body, or social performance. During camp, Aastha and I will check in proactively because some youth will not ask for help. We will give one direction at a time, keep the full sequence visible, make public reading and presenting optional, and let youth pass and re-enter. These supports do not lower the expectation that learners use evidence. They improve the evidence we can gather about their thinking.

Technology will remain optional and purposeful. A phone camera can preserve a transient observation, zoom in on a suspected fiber, or let a learner contribute without copying a large amount of text. A shared slide or projected model can make revisions visible. A timer can support transitions. None of these tools is automatically just because it is digital. I will ask whether a tool increases access to the phenomenon, strengthens evidence, or expands youth agency. If paper, a magnifier, or conversation does the job better, technology should not displace it.

\section{Principle 6: Justice Must Move From Evidence to Collective Action}

The litter investigations connect local evidence with communication and action, while the transdisciplinary and joy-centered readings keep the analysis from becoming a narrow cleanup lesson \citep{smetana2019,wisbrock2020,nesmith2021,muhammad2023}. If we ask only, ``What can an individual camper recycle?'' we hide the roles of product design, production, labor, wastewater systems, policy, corporate decisions, and unequal exposure. If we describe harm without opening a path toward change, we can leave youth with fear or helplessness.

\keyclaim{Justice-centered science must help youth connect mechanism, power, and an action that does not blame the people carrying the greatest burden.} Our model will identify where plastic is produced, released, transported, encountered by living systems, and governed by human choices. Youth advocacy can address product design, reuse and repair, monitoring, indoor or outdoor environments, public communication, research questions, or policy. Art made for a real audience lets youth combine scientific claims with the future they want to argue for.

The action still has to match the evidence. We should not use a body-based anchor to frighten youth or make unsupported health claims. We should not pretend that cleaning one sample solves a system-level problem. Instead, youth can say what the evidence supports, what remains uncertain, who needs better information, and what change would reduce a pathway or distribute responsibility more fairly. That is a stronger form of advocacy because it treats the audience and the young scientist with respect.

\section{My Camp Commitments}

These principles leave me with a practical set of commitments for Invisible Invaders:

\begin{itemize}
  \item \textbf{Plan backward, then revise forward.} I will keep the final explanation and evidence goals clear while changing the route in response to youth questions, observations, access needs, and feedback.
  \item \textbf{Keep the storyline visible.} Every investigation will connect to the driving question and a named Gotta-Have, with the purpose and sequence available for learners who need to pause or re-enter.
  \item \textbf{Protect scientific uncertainty.} Models and public products will distinguish observation from inference, suspected from confirmed particles, and exposure evidence from claims about route or harm.
  \item \textbf{Design equivalent routes into rigor.} Movement and seated options, speaking and non-speaking options, individual thinking and supported collaboration will all lead to evidence-based explanation and action.
  \item \textbf{Track whose thinking becomes public.} I will notice whose questions, model relationships, and proposed actions enter the shared record, especially when a learner contributes quietly or outside the expected form.
  \item \textbf{Use technology for access, evidence, or agency.} Novelty is not a learning goal. A tool earns its place when it makes meaningful scientific participation possible or more powerful.
  \item \textbf{Debrief like a learning scientist.} At the end of each day I will ask what youth evidence changed our plan, which barrier hid someone's thinking, and what Aastha and I need to redesign before the next session.
\end{itemize}

\section{Closing Reflection}

Module 2 did not give me one method for running an inclusive science camp. It gave me a more disciplined way to decide. Coherence asks whether youth can understand and shape the journey; modeling, whether their ideas become public and revisable; place and familial curriculum, whose knowledge enters science; identity and joy, who is recognized as belonging; UDL, PICRAT, and my Puzzle Plan, whether the environment reveals thinking or filters learners through one preferred bodymind; and justice, whether evidence opens power-aware, collective possibility.

\keyclaim{The teacher-as-learning-scientist stance begins by studying the learning environment before explaining the learner as the problem.} In Invisible Invaders, I want that stance to be visible in small choices: one direction at a time, a model that can change, an evidence label that protects uncertainty, an optional public presentation, a proactive check-in, a community question that matters, and more than one credible way to act. Those choices are not separate from strong science teaching. They are how I will know whether our rigorous science became accessible enough for youth to claim as their own.

\section{AI Use Disclosure}

AI was used under Piter Garcia's supervision and guidance to make this document more inclusive and accessible by organizing, structuring, formatting, and enhancing content that Piter provided. Piter guided the analysis, reviewed the course sources and claims, and remains responsible for the final interpretation and submission. AI supported cross-document synthesis, source validation, direct-source linking, and checks for missing causal or access connections; it did not replace Piter's reading, judgment, camp planning, or authorship.

\bibliographystyle{plainnat}
\bibliography{references}

\end{document}
